\chapter{The segmentation head, BCE plus Dice, and auditing labels}\label{ch:seghead}
\begingroup\setlength{\emergencystretch}{5em}% long provenance paths: allow loose lines rather than overfull ones
\def\tbrk{\hfil\penalty50\hfilneg}% local: a line may break before a provenance tag, leaving the line ragged

\begin{onething}
Fine-tuning adds a small segmentation head to the shared encoder: five independent yes/no outputs per pixel (water,
cloud, thin cloud, shadow, snow), each a sigmoid, trained with binary cross-entropy (the negative log-likelihood of a
coin flip) plus soft Dice (an overlap score that a rare class cannot dodge), counted only over pixels whose label is
trusted. Which labels to trust was decided by auditing the provider's labels without any ground truth.
\end{onething}

\section{Why a segmentation head}\label{sec:seghead:why}

\emph{Segmentation}\index{segmentation} means giving every pixel a label, ``water'' or ``cloud'', rather than one label per
image. The product already segments: the \code{scene_segmentation} action (\cref{ch:product}) computes spectral indices
(NDVI for vegetation, NDWI for water, a brightness measure) and applies fixed thresholds to them
(\file{backend/allotrope/action_types/scene_segmentation.py}). Its masks decide which pixels \code{anomaly_scoring} looks at.

Thresholds on indices are blunt. Cloud edges, thin cloud and shadows slip through and get scored as anomalies: a cloud is
something the reconstruction model rarely redraws well, so its error is large, so it looks anomalous. The obvious fix is a
better rule (a lower brightness threshold, a second index), and it fails the same way one notch later: a hand-written rule has
no idea what thin cloud over a dark field looks like until someone writes that case in.

The fine-tuning plan (\file{docs/proposals/fine-tuning-and-segmentation.md}) instead teaches the model. Indradhanu's encoder
has already spent a hundred epochs learning what hyperspectral scenes look like (\cref{ch:pretrain}); a small extra network on
top of its features, a \emph{head}\index{head (network)}, can learn to outline water and cloud from labelled examples at a fraction
of the cost of a new model. The head is judged by detection: nothing ships that lowers the PRISMA bench.

The chapter has two halves. The first (\cref{sec:seghead:arch,sec:seghead:sigmoid,sec:seghead:bce,sec:seghead:dice,sec:seghead:iou,sec:seghead:known})
builds the head and its loss, deriving each formula. The second (\cref{sec:seghead:audit}) asks which of the labels it learns
from can be trusted, when nobody has a correct set to compare with.

\section{The head: the decoder of an earlier chapter, smaller}\label{sec:seghead:arch}

\subsection{What it is}

In \file{app/foundation_models/components/spectral_stem_seg_former_mae.py} the head is one line:
\begin{lstlisting}
self.seg_head = (SegFormerDecoder(embed_dims=list(embed_dims), decoder_dim=seg_decoder_dim,
                                  out_channels=seg_classes) if seg_classes else None)
\end{lstlisting}
It is the same \code{SegFormerDecoder} class that reconstructs the spectrum (\cref{ch:decoder}) \citep{xie2021segformer}, built
with width \code{decoder_dim} $d = 128$ instead of 256 and \code{out_channels} $K = 5$ (one per class) instead of 32. Pretrained
models have \code{seg_classes = 0}, so no head; fine-tuning sets 5.

Recall the decoder's steps. The encoder (\cref{ch:encoder}) produces four feature maps at strides 4, 8, 16 and 32 with 32, 64,
160 and 256 channels. The decoder
\begin{enumerate}
  \item projects each map to $d$ channels with a $1 \times 1$ convolution;
  \item upsamples each to the stride-4 grid (bilinear) and stacks them: $4d$ channels;
  \item fuses them back to $d$ with a $1 \times 1$ convolution and a GELU;
  \item refines with a $3 \times 3$ convolution and a GELU;
  \item predicts $16K$ channels with a second $3 \times 3$ convolution: one value for each of the 16 sub-pixels of a $4 \times 4$
    block, for each of the $K$ classes;
  \item rearranges those channels into a full-resolution map with \emph{PixelShuffle}(4) \citep{shi2016realtime}.
\end{enumerate}

\begin{table}[htbp]
\centering
\caption{Shapes through the segmentation head for a batch of $B$ tiles of 128 $\times$ 128 pixels.}
\label{tab:seghead:shapes}
\begin{tabular}{ll}
\toprule
step & shape\\
\midrule
encoder stage 1, $F_1$ & \shape{(B, 32, 32, 32)}\\
encoder stage 2, $F_2$ & \shape{(B, 64, 16, 16)}\\
encoder stage 3, $F_3$ & \shape{(B, 160, 8, 8)}\\
encoder stage 4, $F_4$ & \shape{(B, 256, 4, 4)}\\
project each to 128, upsample to 32 $\times$ 32, stack & \shape{(B, 512, 32, 32)}\\
fuse ($1 \times 1$) & \shape{(B, 128, 32, 32)}\\
refine 1 ($3 \times 3$) & \shape{(B, 128, 32, 32)}\\
refine 2 ($3 \times 3$), $16 \times 5 = 80$ channels & \shape{(B, 80, 32, 32)}\\
PixelShuffle(4): $80 = 5 \times 4 \times 4$ & \shape{(B, 5, 128, 128)}, logits\\
\bottomrule
\end{tabular}
\end{table}

A \emph{logit}\index{logit} is a raw score, any real number; \cref{sec:seghead:sigmoid} turns it into a probability.

\coursefig[0.9\linewidth]{iw-47-seg-head}{The head in context: the shared stem and encoder feed two heads. Pass 1 removes 65\% of
the tokens and reconstructs; pass 2 keeps every token and segments. Which passes run depends on the training mode
(\cref{ch:finetune}).}{fig:seghead:head}

\begin{sourcenote}[: design text against code]
The proposal's first description of the head was ``1x1 projections to a common width, upsample to 1/4, fuse, 1x1 to classes,
bilinear to full size''. The code reuses \code{SegFormerDecoder}, which ends in two $3 \times 3$ convolutions and a PixelShuffle,
not a $1 \times 1$ and a bilinear upsample (\file{app/foundation_models/components/seg_former_decoder.py}, lines 113--117). The
source wins: the head predicts each pixel of a $4 \times 4$ block separately, as the reconstruction does. The proposal now says
``3x3 refine, PixelShuffle to full size'', so the two agree today.
\end{sourcenote}

\subsection{Counting its parameters before reading them}

A convolution with $C_{\text{in}}$ input channels, $C_{\text{out}}$ output channels and a $k \times k$ kernel has
$k^2 C_{\text{in}} C_{\text{out}}$ weights plus $C_{\text{out}}$ biases (\cref{ch:shapes}). Apply it to every layer, with $d = 128$
and $K = 5$:

\begin{worked}[label={we:seghead:params}]{The head's parameter count}
Do this on paper, one layer per line.
\begin{align*}
  \text{project } 32 \to 128:&\quad 32 \times 128 + 128 = 4{,}096 + 128 = 4{,}224\\
  \text{project } 64 \to 128:&\quad 64 \times 128 + 128 = 8{,}192 + 128 = 8{,}320\\
  \text{project } 160 \to 128:&\quad 160 \times 128 + 128 = 20{,}480 + 128 = 20{,}608\\
  \text{project } 256 \to 128:&\quad 256 \times 128 + 128 = 32{,}768 + 128 = 32{,}896\\
  \text{fuse } 512 \to 128 \ (1\times1):&\quad 512 \times 128 + 128 = 65{,}536 + 128 = 65{,}664\\
  \text{refine 1 } 128 \to 128 \ (3\times3):&\quad 9 \times 128 \times 128 + 128 = 147{,}456 + 128 = 147{,}584\\
  \text{refine 2 } 128 \to 80 \ (3\times3):&\quad 9 \times 128 \times 80 + 80 = 92{,}160 + 80 = 92{,}240.
\end{align*}
Add them in pairs:
\begin{align*}
  4{,}224 + 8{,}320 &= 12{,}544\\
  12{,}544 + 20{,}608 &= 33{,}152\\
  33{,}152 + 32{,}896 &= 66{,}048\\
  66{,}048 + 65{,}664 &= 131{,}712\\
  131{,}712 + 147{,}584 &= 279{,}296\\
  279{,}296 + 92{,}240 &= 371{,}536.
\end{align*}
Building the head and counting its parameters gives 371,536 \tbrk\tagM{indradhanu workings/47-the-segmentation-head.md} exactly. The
GELUs, the bilinear upsampling and the PixelShuffle have no parameters.
\end{worked}

\paragraph{A closed form.} Write the same count with symbols, $d$ for the width and $K$ for the classes. The four input widths sum to
$32 + 64 + 160 + 256 = 512$.
\begin{align*}
  \text{projections:}&\quad (32 + 64 + 160 + 256)\,d + 4d = 512d + 4d\\
  \text{fuse:}&\quad 4d \cdot d + d = 4d^2 + d\\
  \text{refine 1:}&\quad 9d^2 + d\\
  \text{refine 2:}&\quad 9 \cdot d \cdot 16K + 16K = 144Kd + 16K.
\intertext{Collect the $d^2$ terms ($4 + 9 = 13$) and the $d$ terms ($512 + 4 + 1 + 1 = 518$):}
  P(d, K) &= 13d^2 + (518 + 144K)\,d + 16K.
\end{align*}
Check it on the head, $d = 128$, $K = 5$: $13 \times 16{,}384 = 212{,}992$; $(518 + 720) \times 128 = 1{,}238 \times 128 = 158{,}464$;
$16 \times 5 = 80$; total $212{,}992 + 158{,}464 + 80 = 371{,}536$. Check it on the reconstruction decoder, $d = 256$, $K = 32$:
$13 \times 65{,}536 = 851{,}968$; $(518 + 4{,}608) \times 256 = 5{,}126 \times 256 = 1{,}312{,}256$; $16 \times 32 = 512$; total
$851{,}968 + 1{,}312{,}256 + 512 = 2{,}164{,}736$, the figure \cref{ch:params} derived by hand. The head is
$371{,}536 / 2{,}164{,}736 = 0.17$, about a sixth of the reconstruction decoder. The $d^2$ term explains the saving: halving $d$
divides the dominant term by four.

\section{Sigmoid or softmax: one yes/no per class}\label{sec:seghead:sigmoid}

\subsection{The sigmoid as log-odds}

A probability $p$ of ``yes'' can be rewritten as \emph{odds}, $p/(1-p)$: odds of 9 mean ``nine to one on''. The logarithm of the odds,
the \emph{log-odds}, runs over all real numbers: very unlikely events have large negative log-odds, near-certain ones large positive.
So a network that outputs any real number $z$ can be read as log-odds:
\begin{align*}
  z &= \ln\frac{p}{1-p}.
\intertext{Solve for $p$. Exponentiate both sides:}
  e^{z} &= \frac{p}{1-p}.
\intertext{Multiply both sides by $1 - p$ and collect the $p$ terms:}
  e^{z} - e^{z}p &= p \quad\Longrightarrow\quad e^{z} = p\,(1 + e^{z}).
\intertext{Divide by $1 + e^{z}$, then divide numerator and denominator by $e^{z}$:}
  p &= \frac{e^{z}}{1 + e^{z}} = \frac{1}{1 + e^{-z}} = \sigma(z).
\end{align*}
\begin{definition}[Sigmoid]\index{sigmoid}
$\sigma(z) = 1/(1 + e^{-z})$ maps a logit $z$ to a probability in $(0, 1)$; $\sigma(0) = 0.5$, and $p > 0.5$ exactly when $z > 0$.
\end{definition}

\subsection{Multi-label is not multi-class}

\begin{definition}[Multi-class and multi-label]\index{multi-label classification}
A problem is \emph{multi-class} when exactly one of $K$ classes is true for each pixel; it is \emph{multi-label} when any number of
them, from none to all, can be true at once.
\end{definition}
For multi-class problems the standard output is the \emph{softmax}\index{softmax}, which turns $K$ logits into $K$ probabilities that sum
to 1 (\cref{ch:attention}):
\[
  \softmax(\vect{z})_k = \frac{e^{z_k}}{\sum_{j=1}^{K} e^{z_j}}.
\]
The classes here overlap: thin cloud over water is both. So the head uses $K$ independent sigmoids, one per class, with no link between
them. It is five separate yes/no questions asked of the same features.

\begin{worked}[label={we:seghead:softmax}]{Why the softmax cannot say ``water and thin cloud''}
A pixel is thin cloud over water. Give it logits $+2.197$ for water and for thin cloud, and $-2.197$ for cloud, shadow and snow.
Use $e^{2.197} = 9.00$ and $e^{-2.197} = 1/9 = 0.111$ (because $\ln 9 = 2.197$).

Softmax:
\begin{align*}
  \text{sum} &= 9 + 0.111 + 9 + 0.111 + 0.111 = 18.333\\
  \text{water} &= 9 / 18.333 = 0.491\\
  \text{thin cloud} &= 9 / 18.333 = 0.491.
\end{align*}
Each must give up half its probability to the other. Sigmoids:
\begin{align*}
  \text{water} &= 1 / (1 + e^{-2.197}) = 1 / (1 + 0.111) = 1 / 1.111 = 0.900\\
  \text{thin cloud} &= 0.900, \text{ independently.}
\end{align*}
With a softmax the head could never say ``certainly water, and certainly under thin cloud''.
\end{worked}

\paragraph{The link between the two.} A softmax over two logits $(z, 0)$ gives the first class
\[
  \frac{e^{z}}{e^{z} + e^{0}} = \frac{e^{z}}{e^{z} + 1} = \sigma(z).
\]
So one sigmoid is a two-class softmax with the second logit fixed at zero, and five sigmoids are five separate two-class problems
(class or not-class). There is no ``land'' or ``clear'' output: land is water's ``no'' where the pixel is labelled, and clear sky is
cloud's and thin cloud's ``no''.

\subsection{The five classes}

\file{app/utils/segmentation/targets.py} defines the classes and how each is read from EnMAP's quality layers:
\begin{center}\small
\begin{tabular}{clll}
\toprule
index & class & target is 1 where & notes\\
\midrule
0 & water & \texttt{label\_classes == 2} & 1 = land; 3 = no-data\\
1 & cloud & \texttt{label\_cloud == 1} or \texttt{label\_classes == 0} & code 0 is EnMAP's cloud class\\
2 & thin\_cloud & \texttt{label\_cirrus >= 2} or \texttt{label\_haze == 1} & cirrus levels 0--3\\
3 & shadow & \texttt{label\_cloud\_shadow == 1} & \\
4 & snow & \texttt{label\_snow == 1} & \\
\bottomrule
\end{tabular}
\end{center}

\section{Binary cross-entropy from the coin flip}\label{sec:seghead:bce}

\subsection{The likelihood of one label}

Treat one pixel and one class as a coin flip. The head says the probability of ``yes'' is $p$; the label $t$ is 1 (yes) or 0 (no). A
\emph{Bernoulli}\index{Bernoulli distribution} variable with parameter $p$ takes value 1 with probability $p$ and 0 with probability
$1 - p$. Both cases fit in one formula:
\begin{equation}\label{eq:seghead:bern}
  P(t \mid p) = p^{t}\,(1 - p)^{1 - t}.
\end{equation}
Check it: for $t = 1$, $p^1 (1-p)^0 = p \times 1 = p$; for $t = 0$, $p^0 (1-p)^1 = 1 \times (1 - p) = 1 - p$.

\subsection{From likelihood to loss}

The \emph{likelihood}\index{likelihood} of a set of labels is the probability the model gives to all of them at once. Treat the $n$ pixels
as independent, so the probabilities multiply:
\begin{align*}
  \mathcal{P} &= \prod_{i=1}^{n} p_i^{t_i}(1 - p_i)^{1 - t_i}.
\intertext{A good model makes $\mathcal{P}$ large. Products of many numbers below 1 underflow, so take the logarithm, which turns the
product into a sum (because $\ln(ab) = \ln a + \ln b$) and the powers into factors (because $\ln a^{t} = t\ln a$):}
  \ln\mathcal{P} &= \sum_{i=1}^{n}\bigl[t_i \ln p_i + (1 - t_i)\ln(1 - p_i)\bigr].
\intertext{The logarithm is increasing, so maximising $\ln\mathcal{P}$ maximises $\mathcal{P}$. Training minimises, so take the negative,
and divide by $n$ so the loss does not grow with the number of pixels:}
  \mathrm{BCE} &= -\frac{1}{n}\sum_{i=1}^{n}\bigl[t_i \ln p_i + (1 - t_i)\ln(1 - p_i)\bigr].
\end{align*}
\begin{definition}[Binary cross-entropy]\index{binary cross-entropy}
For one pixel, $\ell(p, t) = -\bigl[t\ln p + (1 - t)\ln(1 - p)\bigr]$: $-\ln p$ when the label is 1 and $-\ln(1 - p)$ when it is 0. BCE is
the mean of $\ell$ over pixels: the negative log-likelihood of the labels under the model, per pixel.
\end{definition}

\begin{worked}[label={we:seghead:bce3}]{BCE on three pixels}
Do this on paper with a calculator for the logarithms.
\begin{center}
\begin{tabular}{cccl}
\toprule
pixel & $t$ & $p$ & reading\\
\midrule
1 & 1 & 0.9 & confident and right\\
2 & 0 & 0.5 & unsure\\
3 & 0 & 0.9 & confident and wrong\\
\bottomrule
\end{tabular}
\end{center}
\begin{align*}
  \ell_1 &= -\ln 0.9 = 0.1054\\
  \ell_2 &= -\ln(1 - 0.5) = -\ln 0.5 = 0.6931\\
  \ell_3 &= -\ln(1 - 0.9) = -\ln 0.1 = 2.3026\\
  \text{sum} &= 0.1054 + 0.6931 + 2.3026 = 3.1011\\
  \mathrm{BCE} &= 3.1011 / 3 = 1.0337.
\end{align*}
The confident mistake costs 22 times the confident success ($2.3026 / 0.1054 = 21.8$). As $p \to 1$ on a pixel whose label is 0,
$-\ln(1 - p) \to \infty$: cross-entropy punishes confident errors without limit.
\end{worked}

\subsection{The gradient is $p - t$}

The head outputs logits $z$, and $p = \sigma(z)$. Training needs $\partial\ell/\partial z$. Two steps and the chain rule
(\cref{ch:trainloop}).

\textbf{Step 1: the sigmoid's derivative.} Write $\sigma(z) = (1 + e^{-z})^{-1}$ and differentiate the outer power, then the inside:
\begin{align*}
  \frac{d\sigma}{dz} &= -1 \cdot (1 + e^{-z})^{-2} \cdot \frac{d}{dz}(1 + e^{-z}) = -(1 + e^{-z})^{-2}\cdot(-e^{-z})
                      = \frac{e^{-z}}{(1 + e^{-z})^2}.
\intertext{Split the fraction into two factors:}
  \frac{d\sigma}{dz} &= \frac{1}{1 + e^{-z}}\cdot\frac{e^{-z}}{1 + e^{-z}}.
\intertext{The first factor is $\sigma$. The second is $1 - \sigma$, because $1 - \frac{1}{1 + e^{-z}} = \frac{1 + e^{-z} - 1}{1 + e^{-z}}
= \frac{e^{-z}}{1 + e^{-z}}$. So}
  \frac{d\sigma}{dz} &= \sigma(1 - \sigma) = p(1 - p).
\end{align*}

\textbf{Step 2: the loss's derivative in $p$.} Differentiate $\ell = -t\ln p - (1-t)\ln(1-p)$ term by term (using $\frac{d}{dp}\ln p = 1/p$ and
$\frac{d}{dp}\ln(1-p) = -1/(1-p)$):
\[
  \frac{\partial\ell}{\partial p} = -\frac{t}{p} + \frac{1 - t}{1 - p}.
\]

\textbf{Chain rule.} Multiply the two:
\begin{align*}
  \frac{\partial\ell}{\partial z} &= \Bigl(-\frac{t}{p} + \frac{1 - t}{1 - p}\Bigr)\, p(1 - p)
\intertext{Multiply each term out; the $p$ cancels in the first and the $1 - p$ in the second:}
  &= -t(1 - p) + (1 - t)\,p
\intertext{Expand and cancel $tp$:}
  &= -t + tp + p - tp = p - t.
\end{align*}
\begin{keypoint}
The gradient of binary cross-entropy with respect to the logit is the prediction minus the label, $p - t$: no larger than 1 in size,
and zero only when the prediction equals the label.
\end{keypoint}
For the three pixels of \cref{we:seghead:bce3}: pixel 1, $0.9 - 1 = -0.1$ (push the logit up a little); pixel 2, $0.5 - 0 = 0.5$ (push down);
pixel 3, $0.9 - 0 = 0.9$ (push down hard). The sigmoid's own derivative, $p(1-p)$, becomes tiny as $p$ nears 0 or 1 ($0.99 \times 0.01 = 0.0099$),
yet the loss's gradient does not: the $1/p$ and $1/(1-p)$ of the logarithm cancel it exactly. This is why sigmoid and cross-entropy are always used together.

\paragraph{Numerically safe form.} Computing $\ln\sigma(z)$ for $z = -100$ underflows. PyTorch's
\code{binary_cross_entropy_with_logits}, which \code{masked_bce_dice} calls, takes the logit directly and uses the equivalent form
$\ell = \max(z, 0) - zt + \ln(1 + e^{-|z|})$, which never exponentiates a large positive number. Check it at $z = \ln 9 = 2.197$, $t = 0$:
$\max(2.197, 0) - 0 + \ln(1 + e^{-2.197}) = 2.197 + \ln(1.111) = 2.197 + 0.105 = 2.303 = -\ln 0.1$, as in pixel 3.

\section{Overlap: the Dice coefficient and soft Dice}\label{sec:seghead:dice}

\subsection{Dice for sets}

Let $A$ be the set of pixels the head calls ``cloud'' and $B$ the set the label calls ``cloud''.
\begin{definition}[Dice coefficient]\index{Dice coefficient}
$D(A, B) = \dfrac{2\,\abs{A \cap B}}{\abs{A} + \abs{B}}$: twice the overlap divided by the two sizes added. $D = 1$ for identical sets, $D = 0$
for disjoint ones.
\end{definition}
The factor 2 makes the scale right: if $A = B$ then $\abs{A \cap B} = \abs{A}$ and $D = 2\abs{A}/(2\abs{A}) = 1$.

\subsection{Soft Dice: a differentiable version}

Set membership is all-or-nothing, so its gradient is zero almost everywhere. Replace each pixel's membership of $A$ by its probability
$p_i$ and of $B$ by its label $t_i$ \citep{milletari2016vnet}:
\begin{align*}
  \abs{A \cap B} &\to \sum_i p_i t_i, \qquad \abs{A} \to \sum_i p_i, \qquad \abs{B} \to \sum_i t_i.
\end{align*}
If the $p_i$ are exactly 0 or 1 these are the set sizes again. The code adds 1 to the top and bottom (``smoothing'') so the ratio is
defined when a class is absent from both prediction and label, and turns the coefficient into a loss by subtracting it from 1:
\begin{equation}\label{eq:seghead:dice}
  \mathcal{L}_{\text{Dice}} = 1 - \frac{2\sum_i p_i t_i + 1}{\sum_i p_i + \sum_i t_i + 1}.
\end{equation}
If a class is absent from the label and predicted absent, all sums are 0 and the loss is $1 - 1/1 = 0$: correct, with no division by zero.

\begin{worked}[label={we:seghead:dice2x2}]{Soft Dice on a $2 \times 2$ mask}
Do this on paper. Four pixels, labels and predicted probabilities:
\begin{center}
\begin{tabular}{c|cc}
& left & right\\
\midrule
top & $t = 1$, $p = 0.8$ & $t = 1$, $p = 0.4$\\
bottom & $t = 0$, $p = 0.3$ & $t = 0$, $p = 0.1$\\
\end{tabular}
\end{center}
\begin{align*}
  \textstyle\sum p_i t_i &= 0.8 \times 1 + 0.4 \times 1 + 0.3 \times 0 + 0.1 \times 0 = 0.8 + 0.4 = 1.2\\
  \textstyle\sum p_i &= 0.8 + 0.4 + 0.3 + 0.1 = 1.6\\
  \textstyle\sum t_i &= 1 + 1 + 0 + 0 = 2\\
  \text{soft Dice, no smoothing} &= \frac{2 \times 1.2}{1.6 + 2} = \frac{2.4}{3.6} = 0.667\\
  \text{with smoothing} &= \frac{2.4 + 1}{3.6 + 1} = \frac{3.4}{4.6} = 0.739\\
  \mathcal{L}_{\text{Dice}} &= 1 - 0.739 = 0.261.
\end{align*}
Compare the hard Dice at threshold 0.5. The head calls only the top-left pixel cloud: $A = \{\text{TL}\}$, $B = \{\text{TL}, \text{TR}\}$,
so $\abs{A \cap B} = 1$ and $D = 2 \times 1 / (1 + 2) = 2/3 = 0.667$. Here the soft and hard values agree; in general they do not.
\end{worked}

\subsection{The gradient of soft Dice is global}

Write $I = \sum p_i t_i$ and $S = \sum p_i + \sum t_i$, so the coefficient is $Q = (2I + 1)/(S + 1)$. Differentiate with respect to one
pixel's $p_j$ by the quotient rule ($\partial I/\partial p_j = t_j$, $\partial S / \partial p_j = 1$):
\begin{align*}
  \frac{\partial Q}{\partial p_j} = \frac{2t_j\,(S + 1) - (2I + 1)\cdot 1}{(S + 1)^2}.
\end{align*}
In \cref{we:seghead:dice2x2}, $S + 1 = 4.6$ and $2I + 1 = 3.4$, so $(S+1)^2 = 21.16$:
\begin{align*}
  t_j = 1:&\quad \frac{2 \times 4.6 - 3.4}{21.16} = \frac{9.2 - 3.4}{21.16} = \frac{5.8}{21.16} = 0.274\\
  t_j = 0:&\quad \frac{0 - 3.4}{21.16} = -0.161.
\end{align*}
The loss is $1 - Q$, so its gradient has the opposite sign: raise $p$ on labelled pixels, lower it on the others. Unlike BCE's $p_j - t_j$,
each pixel's Dice gradient depends on $I$ and $S$, sums over the \emph{whole class}. That is what lets Dice see a rare class as a whole.

\section{Why BCE plus Dice: class imbalance with numbers}\label{sec:seghead:imbalance}

\subsection{The cheapest wrong answer for BCE}

Suppose a class covers a fraction $\rho$ of the pixels, and the head ignores the image and predicts one constant probability $p$ everywhere.
The BCE is then
\begin{align*}
  \mathrm{BCE}(p) &= -\bigl[\rho\ln p + (1 - \rho)\ln(1 - p)\bigr],
\intertext{because a share $\rho$ of the terms have $t = 1$ and the rest $t = 0$. Find the best constant: set the derivative to zero,}
  \frac{d\,\mathrm{BCE}}{dp} &= -\frac{\rho}{p} + \frac{1 - \rho}{1 - p} = 0.
\intertext{Multiply both sides by $p(1 - p)$, expand, and cancel $\rho p$:}
  \rho(1 - p) &= (1 - \rho)p \quad\Longrightarrow\quad \rho - \rho p = p - \rho p
  \quad\Longrightarrow\quad p = \rho.
\end{align*}
The best constant is the class's base rate, and its loss is $-[\rho\ln\rho + (1-\rho)\ln(1-\rho)]$.

\begin{worked}[label={we:seghead:imbalance}]{A rare class and a common one under BCE and Dice}
Take 1,000 pixels. Shadow covers 10 of them ($\rho = 0.01$), water 300 ($\rho = 0.3$); illustrative shares\tbrk\tagA\ close to the
training data's 0.09\% shadow and 6.9\% water. Score the constant predictor $p = \rho$.

\textbf{BCE.}
\begin{align*}
  \text{shadow:}\quad & -(0.01\ln 0.01 + 0.99\ln 0.99) = 0.01 \times 4.6052 + 0.99 \times 0.01005\\
                      & = 0.04605 + 0.00995 = 0.0560\\
  \text{water:}\quad & -(0.3\ln 0.3 + 0.7\ln 0.7) = 0.3 \times 1.2040 + 0.7 \times 0.3567\\
                      & = 0.3612 + 0.2497 = 0.6109.
\end{align*}
A head that never finds a single shadow pixel pays only 0.056 on shadow, eleven times less than the same laziness costs on water
($0.6109 / 0.0560 = 10.9$). Averaged over classes, ignoring shadow is nearly free.

\textbf{Dice}, same constant predictor:
\begin{align*}
  \text{shadow:}\quad & \textstyle\sum p t = 10 \times 0.01 = 0.1, \quad \sum p = 1000 \times 0.01 = 10, \quad \sum t = 10\\
                      & \mathcal{L}_{\text{Dice}} = 1 - \frac{2 \times 0.1 + 1}{10 + 10 + 1} = 1 - \frac{1.2}{21} = 1 - 0.057 = 0.943\\
  \text{water:}\quad & \textstyle\sum p t = 300 \times 0.3 = 90, \quad \sum p = 300, \quad \sum t = 300\\
                      & \mathcal{L}_{\text{Dice}} = 1 - \frac{180 + 1}{600 + 1} = 1 - 0.301 = 0.699.
\end{align*}
Dice charges the lazy shadow predictor 0.943 out of a maximum of 1, more than it charges water. A perfect shadow prediction
($\sum pt = 10$, $\sum p = 10$) scores $1 - 21/21 = 0$. Under Dice, missing a rare class costs the whole term.
\end{worked}

So the code adds the two, per class:
\begin{equation}\label{eq:seghead:perclass}
  \mathcal{L}_k = \mathrm{BCE}_k + \mathcal{L}_{\text{Dice},k},
\end{equation}
BCE for well-behaved per-pixel gradients ($p - t$, bounded), Dice so that every class, however rare, carries a full-sized term. The loss
is the mean of $\mathcal{L}_k$ over the classes (\cref{sec:seghead:known} says which classes count).

\section{Measuring it: intersection over union}\label{sec:seghead:iou}

\begin{definition}[IoU]\index{intersection over union}
The \emph{intersection over union} (Jaccard index) of a predicted set $A$ and a target set $B$ is $J = \abs{A \cap B}/\abs{A \cup B}$, with
$\abs{A \cup B} = \abs{A} + \abs{B} - \abs{A \cap B}$. A pixel counts as predicted when $p > 0.5$, which is $z > 0$.
\end{definition}

\paragraph{IoU and Dice measure the same thing on different scales.} Write $I = \abs{A \cap B}$ and $U = \abs{A \cup B} = \abs{A} + \abs{B} - I$.
Then $\abs{A} + \abs{B} = U + I$, and
\begin{align*}
  D &= \frac{2I}{\abs{A} + \abs{B}} = \frac{2I}{U + I}.
\intertext{Divide top and bottom by $U$:}
  D &= \frac{2\,(I/U)}{1 + I/U} = \frac{2J}{1 + J}.
\end{align*}
So $D \ge J$ always, with equality only at 0 and 1. In \cref{we:seghead:dice2x2}, $J = 1/2$ and $D = 2 \times 0.5 / 1.5 = 0.667$.

\begin{worked}[label={we:seghead:iou}]{IoU of two pixel sets, and why the tally pools}
Target water: pixels $\{1, 2, 3, 4\}$; predicted water: $\{2, 3, 4, 5, 6\}$.
\begin{align*}
  \text{intersection} &= \{2, 3, 4\}, \text{ size } 3\\
  \text{union} &= \{1, 2, 3, 4, 5, 6\}, \text{ size } 4 + 5 - 3 = 6\\
  J &= 3 / 6 = 0.50.
\end{align*}
\code{IoUTally} adds intersections and unions over the whole validation set, known pixels only, and divides once at the end. Why not
average each batch's IoU? Suppose batch 1 holds a sliver of cloud, $I = 1$, $U = 2$, and batch 2 a large cloud, $I = 90$, $U = 100$:
\begin{align*}
  \text{mean of ratios} &= (1/2 + 90/100)/2 = (0.5 + 0.9)/2 = 0.70\\
  \text{pooled} &= (1 + 90)/(2 + 100) = 91/102 = 0.892.
\end{align*}
The mean of ratios lets a two-pixel batch count as much as a hundred-pixel one; the pooled ratio weights every pixel equally.
\end{worked}

\section{Labels you cannot fully trust: the known mask}\label{sec:seghead:known}

\subsection{The problem}

EnMAP's labels are the output of the provider's automatic tests, and they are not always right. A clear desert scene checked on
2026-09-30 was labelled 57\% cirrus and 14\% cloud.\tbrk\tagM{indradhanu workings/48-checking-labels.md} Train on that and the head learns
``bare desert equals cloud''. Dropping every doubtful scene would also throw away its water, shadow and snow labels, which may be fine.

So each target carries a second array: a \emph{known mask}\index{known mask}, one per class, saying which pixels' labels count. Unknown
means \emph{no opinion}, not \emph{no}.

\subsection{The rules, from the source}

\code{enmap_targets} in \file{app/utils/segmentation/targets.py} builds targets and known masks, both \shape{(5, H, W)}:
\begin{lstlisting}
# class code 3 = provider no-data: unknown for every class
known = np.repeat((cls != 3)[None], len(CLASSES), axis=0)
# code 0 (EnMAP's cloud class): only "cloud" is known
known[[0, 2, 3, 4]] &= ~cloud_class
# cirrus level 1: unsure, not clear
known[2] &= ~((cirrus > 0) & (cirrus < policy.cirrus_min_level) & ~haze)
# snow wins over cloud
known[1] &= ~snow
# no water truth under cloud
known[0] &= ~(cloud | (cirrus >= policy.cirrus_min_level))
dry = (... elev >= policy.plateau_elevation_m) or (... wv < policy.low_water_vapour)
if dry:
    known[1] = known[2] = False      # cloud, thin_cloud unknown for the whole scene
\end{lstlisting}
In words:
\begin{enumerate}
  \item \textbf{Provider no-data} (class code 3): unknown for every class.
  \item \textbf{Under EnMAP's cloud class} (code 0): cloud is known and true; the other four are unknown.
  \item \textbf{Cirrus level 1} is unknown for thin cloud, neither positive nor negative (unless haze is set).
  \item \textbf{Overlaps}: where snow is set, cloud is unknown; where cloud or cirrus $\ge 2$ is set, water is unknown.
  \item \textbf{Dry scenes}: below 0.75\,g/cm$^2$ of water vapour, cloud and thin cloud are unknown for the whole scene; water, shadow and
    snow still count. An elevation rule (3,500\,m and above) is kept as a guard.
\end{enumerate}
The trainer adds two more: a pixel that is nodata in the cube is unknown for every class (\code{seg_known & mask} in
\code{FinetuneTrainer._seg_batch}), and a tile with less than 40\% valid pixels is left out of the segmentation loss
(\code{MIN_VALID_PIXEL_FRACTION = 0.4}).

\begin{sourcenote}[: the cirrus rule]
The course writes the cirrus rule as \code{known[2] &= ~((cirrus == 1) & ~haze)}. The source writes
\code{(cirrus > 0) & (cirrus < policy.cirrus_min_level) & ~haze} (\file{targets.py}, line 65). With the default
\code{cirrus_min_level = 2} the two are identical, since the only integer level with $0 < c < 2$ is 1; the source form also follows the policy
if the threshold is ever changed.
\end{sourcenote}

\subsection{The masked loss}

Let $k_{ik} \in \{0, 1\}$ say whether pixel $i$'s label for class $k$ is known, and $n_k = \sum_i k_{ik}$. \code{masked_bce_dice} computes, for
each class with $n_k > 0$,
\begin{align}
  \mathrm{BCE}_k &= \frac{1}{n_k}\sum_i k_{ik}\,\ell(p_{ik}, t_{ik}),\label{eq:seghead:mbce}\\
  \mathcal{L}_{\text{Dice},k} &= 1 - \frac{2\sum_i k_{ik}\,p_{ik}\,t_{ik} + 1}{\sum_i k_{ik}\,p_{ik} + \sum_i k_{ik}\,t_{ik} + 1},\label{eq:seghead:mdice}
\end{align}
and returns the mean of $\mathrm{BCE}_k + \mathcal{L}_{\text{Dice},k}$ over the classes that have at least one known pixel in the batch. A class
with no known pixel is \emph{left out} of the average, not counted as zero.

\begin{worked}[label={we:seghead:knownmask}]{The known-mask loss on a $2 \times 2$ tile}
A tile from a dry scene, three classes shown (snow and thin cloud have no known pixel here either). Pixel 4 is provider no-data, so it is
unknown for every class; on a dry scene cloud is unknown everywhere. Allow five minutes.
\begin{center}
\begin{tabular}{llcccc}
\toprule
& & pixel 1 & pixel 2 & pixel 3 & pixel 4\\
\midrule
water & target & 1 & 0 & 0 & --\\
      & $p$ & 0.9 & 0.1 & 0.5 & (0.9, ignored)\\
shadow & target & 0 & 0 & 0 & --\\
      & $p$ & 0.1 & 0.1 & 0.1 & (0.5, ignored)\\
cloud & & \multicolumn{4}{c}{unknown in all four pixels}\\
\bottomrule
\end{tabular}
\end{center}
\textbf{Water}, three known pixels:
\begin{align*}
  \text{BCE terms:}&\quad -\ln 0.9 = 0.1054,\quad -\ln(1 - 0.1) = 0.1054,\quad -\ln(1 - 0.5) = 0.6931\\
  \mathrm{BCE} &= (0.1054 + 0.1054 + 0.6931)/3 = 0.9039 / 3 = 0.3013\\
  \textstyle\sum pt &= 0.9 \times 1 = 0.9, \quad \sum p = 0.9 + 0.1 + 0.5 = 1.5, \quad \sum t = 1\\
  \mathcal{L}_{\text{Dice}} &= 1 - \frac{2 \times 0.9 + 1}{1.5 + 1 + 1} = 1 - \frac{2.8}{3.5} = 1 - 0.8 = 0.2\\
  \text{water term} &= 0.3013 + 0.2 = 0.5013.
\end{align*}
\textbf{Shadow}, three known pixels:
\begin{align*}
  \mathrm{BCE} &= (0.1054 + 0.1054 + 0.1054)/3 = 0.1054\\
  \textstyle\sum pt &= 0, \quad \sum p = 0.1 + 0.1 + 0.1 = 0.3, \quad \sum t = 0\\
  \mathcal{L}_{\text{Dice}} &= 1 - \frac{0 + 1}{0.3 + 0 + 1} = 1 - \frac{1}{1.3} = 1 - 0.7692 = 0.2308\\
  \text{shadow term} &= 0.1054 + 0.2308 = 0.3362.
\end{align*}
\textbf{Cloud} has no known pixel and drops out. The loss is the mean over the two classes present:
\[
  (0.5013 + 0.3362)/2 = 0.8375 / 2 = 0.419.
\]
Calling \code{masked_bce_dice} on these numbers returns 0.4187.\tbrk\tagM{indradhanu workings/47-the-segmentation-head.md} Notice pixel 4: its
water probability of 0.9 would have cost $-\ln 0.1 = 2.30$ had it been treated as ``not water''.
\end{worked}

\coursefig[0.9\linewidth]{iw-47-known-mask}{The worked example. Hatched cells are unknown and contribute nothing. Cloud, unknown on a dry
scene, drops out of the class average; the loss is the mean of water's and shadow's terms.}{fig:seghead:knownmask}

\begin{sourcenote}[: the code's Dice and class weights]
In \code{masked_bce_dice} (\file{targets.py}, lines 76--95) the probabilities are multiplied by the known mask before the sums
(\code{p = torch.sigmoid(logits) * k}), and the denominator is \code{(p + targets * k).sum(...)}: exactly \eqref{eq:seghead:mdice}. Since FT04
an optional \code{class_weights} vector turns the class mean into a weighted mean over the classes present,
$\sum_{k \in \text{present}} w_k\mathcal{L}_k / \sum_{k \in \text{present}} w_k$ (\cref{ch:ftruns}).
\end{sourcenote}

\begin{drift}
A smoke test of the probe stage on one EnMAP shard (dev laptop, 2026-10-16) took water IoU from 0 to 0.16, 0.38 and 0.42 over four small epochs:
the first evidence that F001's frozen features already separate water.\tbrk\tagM{docs/lld/finetune.md} The full probe (FT01) reached water 0.82 and
cloud 0.80 (\cref{ch:ftruns}).
\end{drift}

\section{Checking labels when there is no ground truth}\label{sec:seghead:audit}

\subsection{Why labels need auditing}

The known-mask rules above did not come from a textbook. They come from a label audit run on 2026-10-16 over every EnMAP fine-tuning tile: 104
shards, 15,492 tiles, 210 scenes.\tbrk\tagM{research/experiments/indradhanu-finetune/README.md} The audit overturned the first draft of
\file{targets.py} in three places: it treated class code 0 as no-data, switched cloud labels off above 3,500\,m, and treated cirrus level 1 as
``not thin cloud''. All three were wrong.

EnMAP's labels are \emph{quality layers}: masks the provider's processor computes for every scene with automatic tests (is this pixel bright,
cold, high?). A head trained on them learns their mistakes along with their truths. Nobody has a second, correct set of labels to compare with.
So the audit asks a weaker question that can still be answered: \textbf{does what is labelled ``cloud'' look like cloud?}

\begin{analogy}
You suspect a thermometer is faulty but own no reference thermometer. You can still check it against physics: if it reads 35 degrees while the
pond outside is frozen, the reading is wrong. You cannot say what the right temperature was, but you can find \emph{when} it goes wrong, and if
that is always on frosty mornings, you know which readings to distrust. The audit's plausibility tests are the frozen pond; water vapour is the
frosty morning.
\end{analogy}

\subsection{Test 1: is this ``cloud'' bright and flat?}

A cloud is white, which gives two measurable properties:
\begin{itemize}
  \item \textbf{Bright}: the mean reflectance over 450--650\,nm is above 0.2. Clear land in EnMAP averages 0.07 there.
  \item \textbf{Flat}: take the mean reflectance in four narrow windows, 460--500, 540--580, 640--680 and 840--880\,nm, and divide the smallest by
    the largest. White means roughly equal in all four, so the ratio is at least 0.7.
\end{itemize}
A labelled cloud pixel that passes both is \emph{plausible}.

\begin{worked}[label={we:seghead:plausible}]{Four spectra through the cloud test}
Illustrative values\tbrk\tagA; do each ratio by hand.
\begin{center}\small
\begin{tabular}{lcccccl}
\toprule
& 460--500 & 540--580 & 640--680 & 840--880 & min/max & verdict\\
\midrule
cloud & 0.31 & 0.30 & 0.30 & 0.37 & $0.30/0.37 = 0.81$ & bright, flat\\
vegetation & 0.04 & 0.08 & 0.05 & 0.40 & $0.04/0.40 = 0.10$ & dark\\
bright dry soil & 0.18 & 0.25 & 0.32 & 0.38 & $0.18/0.38 = 0.47$ & bright, not flat\\
snow & 0.55 & 0.54 & 0.52 & 0.48 & $0.48/0.55 = 0.87$ & bright, flat\\
\bottomrule
\end{tabular}
\end{center}
Snow passes, because snow is white too. The audit separates it with the \emph{normalised difference snow index}\index{NDSI}: snow is bright in green
and dark near 1600\,nm, where cloud stays bright.
\begin{align*}
  \text{NDSI} &= \frac{\text{green} - \text{SWIR}_{1550\text{--}1680}}{\text{green} + \text{SWIR}_{1550\text{--}1680}}\\
  \text{snow (SWIR 0.086):}&\quad \frac{0.54 - 0.086}{0.54 + 0.086} = \frac{0.454}{0.626} = 0.73 \quad (> 0.4\text{: snow-like})\\
  \text{cloud (SWIR 0.34):}&\quad \frac{0.30 - 0.34}{0.30 + 0.34} = \frac{-0.04}{0.64} = -0.06 \quad (\text{not snow-like}).
\end{align*}
The SWIR values are the audit's group means: snow 0.086, EnMAP's cloud class 0.339.
\end{worked}

These are plausibility tests, not truth. A thin cloud over dark sea can fail ``bright'' and still be real. The audit uses them to find \emph{where}
failures concentrate, not to relabel pixels.

\subsection{Finding 1: class 0 is the cloud class}

EnMAP's land/water map, \code{label_classes}, has four codes. The shard LLD first described codes 0 and 3 as no-data, and the first
\file{targets.py} followed it. The audit looked at the pixels:
\begin{itemize}
  \item every one of the 1,414,747 class-0 pixels has \code{label_cloud = 1};
  \item they hold 79\% of all cloud-labelled pixels;
  \item their mean visible reflectance is 0.30, against 0.07 for clear land: bright, as cloud is.
\end{itemize}
Class 0 is EnMAP's cloud class; only class 3 is no-data. Treating class 0 as no-data would have thrown away most of the cloud the head is meant to
learn:
\begin{align*}
  \text{training tiles with at least 1\% cloud:}&\quad 755 \text{ counting class 0}, \quad 269 \text{ without}\\
  \text{lost:}&\quad 1 - 269/755 = 1 - 0.356 = 0.644 \approx 64\%.
\end{align*}

\begin{sourcenote}[: the LLDs on class 0 and on cloud\_plausible]
\file{docs/lld/enmap-shards.md} (lines 43--44) now reads ``classes 0 = cloud, 1 land, 2 water, 3 no-data - class 0 was first read as no-data''; the
course's remark that the LLD ``calls it no data'' is out of date, and the data, the source and the LLD now agree. Separately,
\file{docs/lld/finetune.md} lists \code{cloud_plausible(refl, wl)} among the interfaces of \file{app/utils/segmentation/targets.py} (``no caller since
FT04''), but no function of that name exists anywhere under \file{app/} (grep \code{def cloud_plausible}): the audit's test lives in the dev-laptop
scripts, not in the repository. The source wins.
\end{sourcenote}

\subsection{Test 2: is this ``cirrus'' cloud, or the ground?}

\emph{Cirrus} is thin, high ice cloud. A thin cloud adds scattered light to whatever is below it, and that added light shows most where the ground is
dark. Over land the visible is dark (0.07) and the shortwave infrared at 2100--2250\,nm brighter (about 0.16), so real thin cloud brightens the visible
\emph{more} than the SWIR. Inside one tile the audit compares pixels labelled cirrus with clear-land pixels (at least 200 of each). If the cirrus pixels
differ from clear land mainly at 2100--2250\,nm, or hardly at all in the visible (less than 0.005), the label follows the ground (brighter, drier soil),
not a cloud: it is \emph{surface-driven}\index{surface-driven label}.
\begin{center}
\begin{tabular}{llccl}
\toprule
tile & band & clear land & ``cirrus'' & change\\
\midrule
X & visible & 0.070 & 0.074 & $+0.004$\\
  & SWIR & 0.160 & 0.210 & $+0.050$: SWIR rose more, surface-driven\\
Y & visible & 0.070 & 0.110 & $+0.040$\\
  & SWIR & 0.160 & 0.175 & $+0.015$: visible rose more, cloud-like\\
\bottomrule
\end{tabular}
\end{center}

\subsection{Finding 2: dry air, not altitude}

The draft rule switched off cloud labels on scenes above 3,500\,m. No scene in the set reaches it: every scene is in India, and the highest scene-mean
elevation is 2,776\,m. What does predict surface-driven cirrus is the scene's \emph{water vapour}\index{water vapour}: the amount of water in the air
column above each square centimetre of ground.

\paragraph{The unit.} If all of it condensed it would form a layer of liquid water. Liquid water has density 1\,g/cm$^3$, so a column of $m$ grams
over 1\,cm$^2$ is $m$\,cm$^3$ of water, a layer $m$\,cm deep:
\[
  0.75 \text{ g/cm}^2 \div 1 \text{ g/cm}^3 = 0.75 \text{ cm} = 7.5 \text{ mm}.
\]
\begin{center}
\begin{tabular}{lccc}
\toprule
water vapour (g/cm$^2$) & scenes & cirrus 1 surface-driven & cirrus $\ge 2$ surface-driven\\
\midrule
below 0.5 & 3 & 100\% & 100\%\\
0.5--0.75 & 8 & 90\% & 100\%\\
0.75--1.0 & 20 & 79\% & 42\%\\
1.0--1.25 & 33 & 26\% & 10\%\\
1.5--2.0 & 60 & 14\% & 8\%\\
2.0--2.5 & 33 & 66\% & 32\%\\
2.5 and above & 32 & 42\% & 3\%\\
\bottomrule
\end{tabular}
\end{center}
Source: \tbrk\tagM{research/experiments/indradhanu-finetune/README.md}. Below 0.75\,g/cm$^2$, cirrus labels are almost all the ground. Above 1.0
most are not, though not all, and the 2.0--2.5 row shows the relationship is not tidy. A likely physical reason, consistent with the table: cirrus
tests usually rely on a band near 1380\,nm, where water vapour normally absorbs so strongly that the ground is invisible and only high cloud
shows. In very dry air the ground shows through and looks like cirrus. The high scenes that prompted the draft rule were also dry; altitude was
standing in for dryness.

\coursefig[0.95\linewidth]{iw-48-label-qa}{(a) Share of cirrus labels that are surface-driven, by the scene's water vapour, for cirrus level 1 and
level 2 or more; below 0.75\,g/cm$^2$ (dashed) nearly all are. (b) Candidate policies: share of suspect labels each removes against share of
plausible labels it also removes. The adopted rule removes 62.6\% of suspect labels for 13.6\% of plausible ones.}{fig:seghead:labelqa}

\subsection{Choosing the rule by its trade-off}

Any rule that removes bad labels also removes some good ones. The audit scored each candidate on the training split by the share of
\emph{suspect} labels it removes (surface-driven cirrus, implausible cloud) and the share of \emph{plausible} labels it removes too:
\begin{center}\small
\begin{tabular}{lcccc}
\toprule
policy & scenes & tiles hit & suspect removed & plausible removed\\
\midrule
water vapour below 0.50 & 3 & 1.5\% & 44.0\% & 0.2\%\\
\textbf{below 0.75 (adopted)} & \textbf{11} & \textbf{5.6\%} & \textbf{62.6\%} & \textbf{13.6\%}\\
below 0.93 (lowest tenth) & 21 & 10.6\% & 62.6\% & 13.6\%\\
water vapour below 1.0 & 30 & 15.2\% & 66.6\% & 15.7\%\\
elevation above 1,500\,m & 13 & 6.6\% & 43.1\% & 18.8\%\\
elevation above 3,500\,m (draft) & 0 & 0\% & 0\% & 0\%\\
\bottomrule
\end{tabular}
\end{center}
Read the adopted row as a ratio, suspect labels removed per plausible label removed:
\begin{align*}
  \text{adopted:}&\quad 62.6 / 13.6 = 4.6\\
  \text{elevation} > 1{,}500\text{ m:}&\quad 43.1 / 18.8 = 2.3 \quad\text{(half as selective).}
\end{align*}
The lowest-tenth rule (below 0.93) removes exactly the same labels as the fixed 0.75, because the scenes between 0.75 and 0.93 hold no suspect
labels to remove, but it switches cloud labels off on 1,545 tiles instead of 808, nearly twice as many. A fixed value wins. The adopted rule touches
11 training scenes and no test scene (the test split's driest scene has 0.94\,g/cm$^2$).

In the code \code{meta.json} stores \code{scene_wv} in g/cm$^2$ times 1000, so the threshold is 750 (\code{LabelPolicy.low_water_vapour}). A scene with
\code{scene_wv = 698} has $698/1000 = 0.698$\,g/cm$^2 < 0.75$: dry.

\begin{keypoint}
The rule was fixed from the audit before any training and is never tuned on results: a rule tuned on IoU would quietly drop whatever labels the head
found hard, the opposite of an audit.
\end{keypoint}

\subsection{Findings 3 to 5: a cirrus level nobody can read, duplicates, rarity}

\paragraph{Cirrus level 1 is neither yes nor no.} Over water the audit measured the share of pixels brighter than 0.05 at 1550--1680\,nm, where water
is black and cloud is not: clear water 3.9\%, cirrus level 1 1.45\%, cirrus level $\ge 2$ 89\%. Level 1 does not look like cloud, so it cannot be a
positive; it is a label that says ``maybe'', so it cannot be a clean negative. It is unknown for thin cloud.

\paragraph{Snow wins over cloud; no water truth under cloud.} 122,815 pixels are labelled both snow and cloud; snow is white, so the cloud test fires
on it, and cloud is made unknown there. Pixels labelled water \emph{and} cirrus $\ge 2$ have a mean near-infrared reflectance of 0.179, against 0.024
for clear water: land-bright, not water. The water label there probably comes from a map of where water usually is, so water is unknown under cloud.

\paragraph{Four scenes counted twice.} Four scenes appear under two processor versions (576 tiles, all in training). A skip list,
\file{research/experiments/indradhanu-finetune/band_lists/enmap_drop_duplicate_scenes.json}, drops them in the loader workers by scene id.

\paragraph{Some classes are rare.} After the policy, in the training split: water 6.9\% of labelled pixels, thin cloud 1.97\%, cloud 0.66\%, snow
0.094\%, shadow 0.086\%. Shadow is about one pixel in $100 / 0.086 = 1{,}163$. Snow has one tile in the test split (639 pixels), too few to measure an
IoU at all. And 53\% of shadow pixels have no cloud within 900\,m in the same tile: an upper bound, since the cloud may sit in the neighbouring tile,
but a reason to treat shadow labels with care too.

\begin{confusions}
\confusion{Is there a ``clear'' or ``land'' output?}{No. Five classes, each yes/no. Land is water's ``no'' where the pixel is labelled; clear sky is
cloud's and thin cloud's ``no''.}
\confusion{Unknown pixels are counted as negatives, then?}{No; that is the whole point of the known mask. They contribute nothing to BCE, to Dice or to
IoU, and a class with no known pixel leaves the class average.}
\confusion{Why not use only Dice, since it handles rare classes?}{Dice's gradient for each pixel depends on whole-class sums and is unstable when
the sums are small; BCE gives every pixel a bounded, local gradient $p - t$. The two together give stable training and full weight to rare classes.}
\confusion{Does adding the head change reconstruction?}{Not by itself: an untrained head leaves the reconstruction pass untouched (a test pins this).
Training the encoder on the segmentation loss can, which is why every stage is gated on the detection bench (\cref{ch:finetune}).}
\confusion{The audit says these labels are wrong; why not relabel them?}{The tests are plausibility checks, not truth. They find \emph{where} failures
concentrate; they cannot say what each pixel really is. ``Unknown'' is the honest answer.}
\confusion{Does the dry-air rule throw away those scenes?}{No. Only cloud and thin cloud become unknown; their water, shadow and snow labels still teach.}
\confusion{Why a whole-scene rule rather than a per-pixel one?}{Water vapour is a scene property and the failures follow it. Per-pixel tests would bake
the audit's own heuristics into the training targets.}
\end{confusions}

\begin{checkyourself}
\item\label{ex:seghead:1} Use the closed form $P(d, K) = 13d^2 + (518 + 144K)d + 16K$ to compute the head's parameter count for $d = 64$, $K = 5$, and
  check it layer by layer.
\item\label{ex:seghead:2} A pixel has logits $+1.0$ for water and thin cloud and $-3.0$ for the other three. What does the sigmoid give for water
  ($e^{-1} = 0.368$)? What would a softmax over the five logits give for water ($e^{-3} = 0.0498$, $e^{1} = 2.718$)?
\item\label{ex:seghead:3} Compute the BCE and its gradient $p - t$ for $t = 1$, $p = 0.2$ and for $t = 0$, $p = 0.2$. Then verify the gradient for the
  first pixel numerically: its logit is $z = \ln(0.2/0.8) = -1.386$; compute $\ell$ at $z = -1.376$ and at $z = -1.396$ and divide the difference by 0.02.
\item\label{ex:seghead:4} Redo the water term of \cref{we:seghead:knownmask} with pixel 3's probability at 0.2 instead of 0.5. Then compute the soft
  Dice loss of \cref{we:seghead:dice2x2} if the top-right probability rises from 0.4 to 0.9.
\item\label{ex:seghead:5} Target cloud covers 8 pixels, the prediction 10 pixels, 6 of them correct. Compute IoU and Dice, and check $D = 2J/(1+J)$.
\item\label{ex:seghead:6} A labelled cloud pixel has window means 0.22, 0.24, 0.25, 0.36 and visible mean 0.24. Is it plausible? A scene has
  \code{scene_wv = 812}: are its cloud labels used? Which ratio of suspect to plausible labels does ``water vapour below 1.0'' achieve, and why was it
  not chosen?
\end{checkyourself}

\nextchapter{ch:finetune}
\endgroup
